% answers_short.tex —— 1977 年·江苏镇江地区初试卷 · 简短答案速查
% 由 tools/make_short_answers.py 自 answers.tex 的 \gksolhead/\gkanswer 提取，
% 并与 1977/1977_数学真题_江苏镇江地区初试_答案要点.md 逐题比对后生成，勿手改（改源 MD 后重跑）。

\begin{gknotice}
\textbf{使用说明}：本节只给最终答案便于快速核对，完整推导见解析版。答案为编写者独立推导，\textbf{非官方标准答案}。
\end{gknotice}

\gkpart{题目}
\begin{itemize}[leftmargin=2.6em]
  \item[\textbf{1.}] (1) $-1982$；(2) $-1$；(3) $\dfrac{15}{2}$；(4) $16$
  \item[\textbf{2.}] (1) $\dfrac{\sqrt2}{4}$；(2) 见解析版
  \item[\textbf{3.}] $x=5$。
  \item[\textbf{4.}] (1) 全体实数（$x\in\mathbb R$）；(2) $x\ne-3$ 且 $x\ne1$；(3) $x\le5$；(4) $1<x<4$
  \item[\textbf{5.}] 见解析版
  \item[\textbf{6.}] $AB=20\sqrt2\ \mathrm m\approx28.3\ \mathrm m$。
  \item[\textbf{7.}] (1) 顶点 $(0,2)$，对称轴 $x=0$，当 $x=0$ 时 $y$ 取极小值 $2$；(2) 新坐标系中方程为 $x_1^2=y_1$，准线 $y_1=-\dfrac14$，焦点 $\left(0,\dfrac14\right)$
  \item[\textbf{8.}] 第一年 $400$ 吨，第二年 $500$ 吨，第三年 $600$ 吨，第四年 $720$ 吨。
  \item[\textbf{9.}] $\dfrac{AG}{GE}=7$。
\end{itemize}
